\begin{lemma}[Affine normalization of $p$-harmonic functions]
Let $1<p<2$, let $\Omega\subseteq\mathrm{Plane}$, let $u\colon\mathrm{Plane}\to\mathbb{R}$ and
$g\colon\mathrm{Plane}\to\mathrm{Plane}$, let $x_0\in\mathrm{Plane}$ and let $r,M,c$ be real
numbers with $r>0$ and $M>0$.  Assume that $u$ is $p$-harmonic in $\Omega$ with weak gradient
$g$.  Put
\[
  T(x):=x_0+rx,\qquad
  \Omega':=\{x\in\mathrm{Plane}:T(x)\in\Omega\},
\]
\[
  v(x):=\frac{u(T(x))-c}{M},\qquad
  G(x):=\frac rM\,g(T(x)).
\]
Then $v$ is $p$-harmonic in $\Omega'$ with weak gradient $G$.
\end{lemma}
\begin{proof}
Throughout, $T(x)=x_0+rx$ is an affine homeomorphism of $\mathrm{Plane}$ with inverse
$T^{-1}(y)=r^{-1}(y-x_0)$, and $\Omega'=T^{-1}(\Omega)$, so $T(\Omega')=\Omega$ and
$T^{-1}(\Omega)=\Omega'$.  Note $r/M>0$.  We verify in turn the six conditions in the
definition of a $p$-harmonic function with weak gradient, \noderef{IsPHarm}.

\medskip\noindent\emph{(1) $\Omega'$ is open.}
By condition (1) of \noderef{IsPHarm}, $\Omega$ is open.  The map $T(x)=x_0+rx$ is continuous,
being the sum of the constant map $x\mapsto x_0$ and the scalar multiplication $x\mapsto rx$.
Hence $\Omega'=T^{-1}(\Omega)$ is the preimage of an open set under a continuous map, so
$\Omega'$ is open.  In particular $\Omega$ and $\Omega'$ are both Lebesgue measurable, which
is used below.

\medskip\noindent\emph{(2) $v$ is continuous on $\Omega'$.}
The affine map $T$ is continuous and maps $\Omega'$ into $\Omega$, and $u$ is continuous on
$\Omega$ by condition (2) of \noderef{IsPHarm}.  Hence $u\circ T$ is continuous on $\Omega'$,
and so is $v=(u\circ T-c)/M$, being obtained from $u\circ T$ by subtracting the constant $c$
and multiplying by the constant $1/M$.

\medskip\noindent\emph{(3) $G$ is almost everywhere strongly measurable on $\Omega'$.}
By condition (3) of \noderef{IsPHarm} there are a strongly measurable
$\tilde g\colon\mathrm{Plane}\to\mathrm{Plane}$ and a Lebesgue-null set
$N\subseteq\mathrm{Plane}$ with $g=\tilde g$ on $\Omega\setminus N$.  The map $T$ is
continuous, hence Borel measurable, so $\tilde g\circ T$ is strongly measurable.  By the
affine change-of-variables background proposition, $T^{-1}$ maps null sets to null sets, so
$T^{-1}(N)$ is null; and $g\circ T=\tilde g\circ T$ on
$\Omega'\setminus T^{-1}(N)=T^{-1}(\Omega\setminus N)$.  Thus $g\circ T$ agrees almost
everywhere on $\Omega'$ with the strongly measurable map $\tilde g\circ T$, and the same holds
for its constant multiple $G=(r/M)(g\circ T)$.

\medskip\noindent\emph{(4) $\lvert G\rvert^{p}$ is integrable on every compact
$K\subseteq\Omega'$.}
Fix such a $K$.  Then $T(K)$ is compact, being the image of a compact set under the continuous
map $T$, and $T(K)\subseteq T(\Omega')=\Omega$.  Since $r/M>0$, for every $x$
\[
  \lvert G(x)\rvert^{p}
  =\Bigl(\tfrac rM\,\lvert g(T(x))\rvert\Bigr)^{p}
  =\Bigl(\tfrac rM\Bigr)^{p}\lvert g(T(x))\rvert^{p},
\]
using that real powers of nonnegative reals are multiplicative,
$(ab)^{p}=a^{p}b^{p}$ for $a,b\ge0$.  By condition (4) of \noderef{IsPHarm} the function
$y\mapsto\lvert g(y)\rvert^{p}$ is integrable on the compact set $T(K)\subseteq\Omega$.  Since
$T$ is injective, $T^{-1}(T(K))=K$, so the affine change-of-variables background proposition,
applied with $S=T(K)$ and $f(y)=\lvert g(y)\rvert^{p}$, shows that
$x\mapsto\lvert g(T(x))\rvert^{p}$ is integrable on $K$ with
\[
  \int_{K}\lvert g(T(x))\rvert^{p}\,dx=r^{-2}\int_{T(K)}\lvert g(y)\rvert^{p}\,dy<\infty .
\]
Multiplying by the finite constant $(r/M)^{p}$ gives the integrability of
$\lvert G\rvert^{p}$ on $K$.

\medskip\noindent\emph{Setting up the test-function correspondence.}
Let $\varphi\colon\mathrm{Plane}\to\mathbb{R}$ be of class $C^\infty$ with compact support and
$\operatorname{supp}\varphi\subseteq\Omega'$.  Define
\[
  \psi:=\varphi\circ T^{-1},\qquad\text{that is}\qquad \psi(y)=\varphi\bigl(r^{-1}(y-x_0)\bigr).
\]
Then:
\begin{enumerate}
\item[(i)--(ii)] $\psi$ is of class $C^\infty$, has compact support, and satisfies
$\operatorname{supp}\psi\subseteq\Omega$, so it is an admissible test function for $\Omega$.
These are clauses (1), (2) and (3) of \noderef{AffineTestFunctionPullback}, applied with this
$x_0$, this $r>0$, the set $\Omega$ and the test function $\varphi$ --- admissible for
$\Omega'=T^{-1}\Omega$ by hypothesis, which is that node's hypothesis on
$\operatorname{supp}\varphi$.  That node needs nothing of $\Omega$ beyond its being a set.
\item[(iii)] $\psi(T(x))=\varphi(x)$ for every $x$, since $T^{-1}(T(x))=x$.
\item[(iv)] Clause (4) of \noderef{AffineTestFunctionPullback} is the chain rule for this
substitution, $D\varphi(x)=r\,D\psi(T(x))$; its inner map $T^{-1}$ is affine with derivative
the scaling by $r^{-1}$, which is why the factor on the right is $r$ and not $r^{-1}$.  So
\begin{equation}\label{eq:phn-chain}
  D\varphi(x)=r\,D\psi(T(x))\qquad\text{for every }x\in\mathrm{Plane}.
\end{equation}
\item[(v)] $D\psi$ vanishes outside $\operatorname{supp}\psi$: indeed
$\mathrm{Plane}\setminus\operatorname{supp}\psi$ is open and $\psi\equiv0$ there, so all
derivatives of $\psi$ vanish on it.  The same remark applies to $\varphi$ and $D\varphi$.
\end{enumerate}
Write $K_0:=\operatorname{supp}\psi$, a compact subset of $\Omega$, and let
$B:=\sup_{y\in\mathrm{Plane}}\lvert D\psi(y)\rvert<\infty$, finite because $D\psi$ is
continuous and vanishes off the compact set $K_0$.

\medskip\noindent\emph{Integrability of the four integrands over $\Omega$.}
We record once and for all that the following functions are integrable on $\Omega$, each
being measurable and vanishing off $K_0$ by (v):
\[
  y\mapsto D\psi(y),\qquad
  y\mapsto u(y)\,D\psi(y),\qquad
  y\mapsto \psi(y)\,g(y),\qquad
  y\mapsto \bigl\langle A_p(g(y)),D\psi(y)\bigr\rangle .
\]
The set $\Omega$ is measurable, being open by condition (1) of \noderef{IsPHarm}; this is what
licenses every measurability claim that follows, since continuity on a set yields almost
everywhere strong measurability with respect to the restricted measure only when that set is
measurable.

For the first two functions, $D\psi$ is continuous on all of $\mathrm{Plane}$ and $u$ is
continuous on $\Omega$ by condition (2) of \noderef{IsPHarm}, so both $D\psi$ and
$u\,D\psi$ are continuous on the measurable set $\Omega$ and hence almost everywhere strongly
measurable with respect to Lebesgue measure restricted to $\Omega$.  Both vanish off the
compact set $K_0\subseteq\Omega$ by (v), and a function continuous on $\Omega$ is bounded on
the compact subset $K_0$, hence integrable there; so each is integrable on $\Omega$ and its
integral over $\Omega$ reduces to its integral over $K_0$.  For the last two we use the
pointwise elementary inequality
\begin{equation}\label{eq:phn-elem}
  t^{a}\le 1+t^{p}\qquad\text{for all }t\ge0\text{ and }0\le a\le p,
\end{equation}
valid because $t^{a}\le1$ when $t\le1$ and $t^{a}\le t^{p}$ when $t\ge1$.  Taking $a=1$ in
\eqref{eq:phn-elem} gives $\lvert\psi(y)g(y)\rvert\le B'\bigl(1+\lvert g(y)\rvert^{p}\bigr)$
on $K_0$ with $B':=\sup\lvert\psi\rvert<\infty$, and $1+\lvert g\rvert^{p}$ is integrable on
$K_0$ by condition (4) of \noderef{IsPHarm} together with $\lvert K_0\rvert<\infty$; combined
with the measurability of $g$ from condition (3) of \noderef{IsPHarm} this gives the
integrability of $\psi g$ on $\Omega$.  For the last function, note from the definition of the
$p$-flux \noderef{PFlux} that $\lvert A_p(\xi)\rvert=\lvert\xi\rvert^{p-2}\lvert\xi\rvert
=\lvert\xi\rvert^{p-1}$ when $\xi\neq0$, while $A_p(0)=0$ and also $0^{p-1}=0$ because
$p-1\neq0$; hence $\lvert A_p(\xi)\rvert=\lvert\xi\rvert^{p-1}$ for every $\xi$.  By
Cauchy--Schwarz and \eqref{eq:phn-elem} with $a=p-1\in[0,p]$,
\[
  \bigl\lvert\bigl\langle A_p(g(y)),D\psi(y)\bigr\rangle\bigr\rvert
  \le B\,\lvert g(y)\rvert^{p-1}\le B\bigl(1+\lvert g(y)\rvert^{p}\bigr),
\]
which is integrable on $K_0$ as before.  The same four estimates, composed with $T$ and
multiplied by constants, give the integrability on $\Omega'$ of the corresponding integrands
there, via the affine change-of-variables background proposition, whose hypothesis that the
set be measurable is met because $\Omega$ and $\Omega'$ are open by step (1).

\medskip\noindent\emph{(5) $G$ is the distributional gradient of $v$ on $\Omega'$.}
This is \noderef{NormalizedWeakGradient}, applied with the same $p$, $\Omega$, $u$, $g$, $x_0$,
$r$, $M$ and $c$, and with the hypothesis that $u$ is $p$-harmonic on $\Omega$ with weak
gradient $g$, which is the hypothesis of the present statement.  Its conclusion is condition
(5) of \noderef{IsPHarm} for $v$ and $G$ on $\Omega'$, verbatim.

That node needs no hypothesis $M>0$: the normalizing factor enters linearly on both sides, so
the identity holds for every $M$.  It also discharges its own inputs rather than taking them
from here --- it applies \noderef{AffineTestFunctionPullback} to push a test function for
$\Omega'$ forward to one for $\Omega$, together with the chain rule \eqref{eq:phn-chain}, and
\noderef{TestFunctionGradientIntegral} to kill the term carrying the additive constant $c$,
which is why $c$ does not appear in $G$.  The two Jacobian factors from the affine change of
variables occur on the two sides of its computation and cancel, so no dimensional constant
enters.

\medskip\noindent\emph{(6) $v$ solves the weak $p$-Laplace equation on $\Omega'$.}
First we compute the flux of $G$.  Fix $x\in\mathrm{Plane}$ and write $\xi:=g(T(x))$, so
$G(x)=(r/M)\xi$.  If $\xi=0$ then $G(x)=0$ and, by \noderef{PFlux}, $A_p(G(x))=0=A_p(\xi)$,
so the identity
\begin{equation}\label{eq:phn-flux}
  A_p(G(x))=\Bigl(\frac rM\Bigr)^{p-1}A_p\bigl(g(T(x))\bigr)
\end{equation}
holds trivially.  If $\xi\neq0$ then $\lvert G(x)\rvert=(r/M)\lvert\xi\rvert>0$ and, by
multiplicativity of real powers on positive reals,
\[
  A_p(G(x))
  =\lvert G(x)\rvert^{p-2}G(x)
  =\Bigl(\frac rM\Bigr)^{p-2}\lvert\xi\rvert^{p-2}\cdot\frac rM\,\xi
  =\Bigl(\frac rM\Bigr)^{p-1}\lvert\xi\rvert^{p-2}\xi
  =\Bigl(\frac rM\Bigr)^{p-1}A_p(\xi),
\]
which is \eqref{eq:phn-flux}.  Now, by \eqref{eq:phn-flux}, \eqref{eq:phn-chain} and
bilinearity of the inner product,
\[
  \int_{\Omega'}\bigl\langle A_p(G(x)),D\varphi(x)\bigr\rangle\,dx
  =\Bigl(\frac rM\Bigr)^{p-1}r\int_{\Omega'}J(T(x))\,dx,
  \qquad J(y):=\bigl\langle A_p(g(y)),D\psi(y)\bigr\rangle ,
\]
and $J$ is integrable on $\Omega$ by the integrability list.  The affine change-of-variables
background proposition with $S=\Omega$ gives
\[
  \int_{\Omega'}J(T(x))\,dx=r^{-2}\int_\Omega J(y)\,dy ,
\]
and $\int_\Omega J(y)\,dy=0$ by condition (6) of \noderef{IsPHarm} applied to the admissible
test function $\psi$.  Hence
\[
  \int_{\Omega'}\bigl\langle A_p(G(x)),D\varphi(x)\bigr\rangle\,dx
  =\Bigl(\frac rM\Bigr)^{p-1}r^{-1}\cdot0=0 ,
\]
which is condition (6) for $v$ and $G$ on $\Omega'$.

All six conditions hold, so $v$ is $p$-harmonic in $\Omega'$ with weak gradient $G$.
\end{proof}
